% Options for packages loaded elsewhere
\PassOptionsToPackage{unicode}{hyperref}
\PassOptionsToPackage{hyphens}{url}
\PassOptionsToPackage{dvipsnames,svgnames,x11names}{xcolor}
\documentclass[
  11pt,
]{article}
\usepackage{fontspec} % arXiv AutoTeX -> XeLaTeX
\usepackage{xcolor}
\setlength{\emergencystretch}{3em}
\usepackage[margin=1in]{geometry}
\usepackage{amsmath,amssymb}
\setcounter{secnumdepth}{-\maxdimen} % remove section numbering
\usepackage{iftex}
\ifPDFTeX
  \usepackage[T1]{fontenc}
  \usepackage[utf8]{inputenc}
  \usepackage{textcomp} % provide euro and other symbols
\else % if luatex or xetex
  \usepackage{unicode-math} % this also loads fontspec
  \defaultfontfeatures{Scale=MatchLowercase}
  \defaultfontfeatures[\rmfamily]{Ligatures=TeX,Scale=1}
\fi
\usepackage{lmodern}
\ifPDFTeX\else
  % xetex/luatex font selection
\fi
% Use upquote if available, for straight quotes in verbatim environments
\IfFileExists{upquote.sty}{\usepackage{upquote}}{}
\IfFileExists{microtype.sty}{% use microtype if available
  \usepackage[]{microtype}
  \UseMicrotypeSet[protrusion]{basicmath} % disable protrusion for tt fonts
}{}
\makeatletter
\@ifundefined{KOMAClassName}{% if non-KOMA class
  \IfFileExists{parskip.sty}{%
    \usepackage{parskip}
  }{% else
    \setlength{\parindent}{0pt}
    \setlength{\parskip}{6pt plus 2pt minus 1pt}}
}{% if KOMA class
  \KOMAoptions{parskip=half}}
\makeatother
\usepackage{longtable,booktabs,array}
\newcounter{none} % for unnumbered tables
\usepackage{calc} % for calculating minipage widths
% Correct order of tables after \paragraph or \subparagraph
\usepackage{etoolbox}
\makeatletter
\patchcmd\longtable{\par}{\if@noskipsec\mbox{}\fi\par}{}{}
\makeatother
% Allow footnotes in longtable head/foot
\IfFileExists{footnotehyper.sty}{\usepackage{footnotehyper}}{\usepackage{footnote}}
\makesavenoteenv{longtable}
\usepackage{graphicx}
\makeatletter
\newsavebox\pandoc@box
\newcommand*\pandocbounded[1]{% scales image to fit in text height/width
  \sbox\pandoc@box{#1}%
  \Gscale@div\@tempa{\textheight}{\dimexpr\ht\pandoc@box+\dp\pandoc@box\relax}%
  \Gscale@div\@tempb{\linewidth}{\wd\pandoc@box}%
  \ifdim\@tempb\p@<\@tempa\p@\let\@tempa\@tempb\fi% select the smaller of both
  \ifdim\@tempa\p@<\p@\scalebox{\@tempa}{\usebox\pandoc@box}%
  \else\usebox{\pandoc@box}%
  \fi%
}
% Set default figure placement to htbp
\def\fps@figure{htbp}
\makeatother
\setlength{\emergencystretch}{3em} % prevent overfull lines
\providecommand{\tightlist}{%
  \setlength{\itemsep}{0pt}\setlength{\parskip}{0pt}}
\usepackage{bookmark}
\IfFileExists{xurl.sty}{\usepackage{xurl}}{} % add URL line breaks if available
\urlstyle{same}
\hypersetup{
  colorlinks=true,
  linkcolor={black},
  filecolor={Maroon},
  citecolor={Blue},
  urlcolor={blue},
  pdfcreator={LaTeX via pandoc}}

\author{}
\date{}

\begin{document}

\section{Absorbing Shadow Classes in Conjunctive Label
Rules}\label{absorbing-shadow-classes-in-conjunctive-label-rules}

\subsubsection{A control-class contamination in a recursive-recommender
benchmark, a protocol that detects it, and a pre-registered null on the
repaired
instrument}\label{a-control-class-contamination-in-a-recursive-recommender-benchmark-a-protocol-that-detects-it-and-a-pre-registered-null-on-the-repaired-instrument}

\textbf{David Mullett} \emph{Independent researcher} · ORCID
0009-0004-2543-1664 · correspondence: d@loopzero.org

\subsection{Abstract}\label{abstract}

Benchmarks for early-warning detection in recursive systems label each
unit as a degradation case or a control, then ask whether a detector
fires before degradation. We show that a natural construction for such a
label --- declaring degradation only when a run condition co-occurs with
an \emph{eligibility} condition --- can silently corrupt the control
class, and we report what happens when the corruption is repaired.

In a MovieLens-25M recommender replay, the benchmark's label rule
declares degradation only when eight consecutive frontier misses are
accompanied by a remaining frontier of at least ten items. The
eligibility clause is absorbing: once a unit's frontier falls below the
floor it can never satisfy the rule again, so units that fail
continuously are filed as controls. In the landmark population, 69.2\%
of controls are frontier-starved and no events are; the ``controls''
miss a median of 15 of 20 pre-landmark steps against at most 7 for
events, making the most natural early-warning statistic
anti-informative. A pre-registered learnability gate detected this
before any detector result existed. We prove the contamination is
structural: because starvation is absorbing, the clean-control set is
invariant to the landmark step, so no landmark choice rescues it.

Treating starvation as a competing terminal event and restricting the
control class accordingly, we then execute a pre-registered evaluation
on the purified population. (A conformance audit before execution found
four departures of the executed configuration from the registration; an
earlier execution, which had halted at its own gate without producing a
detector comparison, was vacated and the registered design executed
once. All findings are reported.) The gate, which failed before repair,
passes after it (TPR 0.561 against a chance rate of 0.05) --- the first
valid measurement this benchmark has produced. On that instrument the
pre-registered slate-structure composite fails all three evaluated
criteria: it is \emph{worse} than the best miss-based baseline at every
matched budget (ΔTPR = −0.472, CI {[}−0.496, −0.444{]} at the primary),
and two of its three covariate coefficients are inverted while the set
carries substantial incremental discrimination (0.162, CI {[}0.143,
0.183{]}).

Purification eliminates or reverses the large contaminated-population
effects: churn inverts (+2.01 \(\rightarrow\) −0.64), occupancy
collapses to null (+0.05), and the outcome-stream quantity reverses
\emph{into} the informative direction (−2.62 \(\rightarrow\) +0.81).
Coverage alone retains a moderate effect in the hypothesised direction
(−0.68): on marginal effects one indicator is inverted, one is null, and
one is consistent with the theory. The composite separates pre-landmark
starved units from clean controls at d = +4.20: what the contaminated
benchmark rewarded was frontier depletion, not early warning.

The failure mode is not new: it is an absorbing competing-risk
pathology. We repair the binary benchmark by landmarking the at-risk
population and excluding subsequent competing-event units from the
residual control class. The contribution is its unrecognized instance in
algorithmically-labeled benchmark construction, and a detection and
audit protocol (a learnability gate plus a shadow-class census) that
exposed both the original label pathology and a subsequent nonconformant
execution before either could support a detector conclusion.

\subsection{Part 1 --- The benchmark and the claim it was built to
test}\label{part-1-the-benchmark-and-the-claim-it-was-built-to-test}

Recursive systems --- those whose current output conditions their future
input --- recur across recommender systems, model-training loops, and
market microstructure. A natural question is whether such systems emit
an early-warning signal before they lose the ability to make progress
--- a signal whose detectability has known limits even in
well-characterised dynamical systems (Boettiger and Hastings, 2012).
Answering it requires a benchmark: a population of units, each labeled
as a degradation case or a control by an operational rule, over which
candidate detectors can be compared at matched false-positive budgets.

We study one such benchmark: a deterministic offline replay on
MovieLens-25M (Harper and Konstan, 2015). Each user is replayed for up
to 50 recursive steps through a frozen item-item collaborative-filtering
engine. At each step the engine recommends a top-K slate; a ``hit''
recovers a held-out positive item and appends it to the profile; the
user's \emph{frontier} is the set of remaining recoverable positives. A
unit is labeled \textbf{degraded} when it accrues a run of eight
consecutive frontier misses while its remaining frontier is at least
ten, within the horizon; otherwise it is a \textbf{control}. A prior
preprint (arXiv:2606.00329) reported a directional early-warning
association on this benchmark using three slate-derived indicators.

Two clarifications frame what follows. First, that prior study's
specific indicators were subsequently found to be partly coupled to the
label rule, which motivated a redefinition. The problem we describe here
is deeper and independent of any indicator: it is a defect in how the
benchmark constructs its \emph{labels}, and it corrupts the target
against which any detector is evaluated. Second, and this governs the
paper's contribution: \textbf{the failure mode is not a new phenomenon.}
It is a concrete instance of hazards that are well developed elsewhere:
informative censoring and competing risks in survival analysis,
depletion of susceptibles and immortal-time bias in
pharmacoepidemiology, negative-class contamination in positive-unlabeled
learning, and reject inference in credit scoring. Our claim is not
phenomenon-novelty. It is that this hazard has an unrecognized and
undetected instance in the construction of recursive-recommender
benchmarks, that its presence inverts the sign of the most natural
early-warning signal, and that a lightweight, pre-registerable protocol
detects it before any result is produced. Throughout we use the terms
\emph{degradation} and \emph{frontier failure} rather than
\emph{collapse}.

\subsection{Part 2 --- The pathology, and the gate that caught it before
any result
existed}\label{part-2-the-pathology-and-the-gate-that-caught-it-before-any-result-existed}

\subsubsection{2.1 The eligibility conjunct and the absorbing shadow
state}\label{the-eligibility-conjunct-and-the-absorbing-shadow-state}

The degradation rule is a conjunction. Let \(m_t\) be the count of
consecutive frontier misses at step t and \(f_t\) the remaining frontier
size. Degradation is declared at the first step where

\(m_t\) \(\geq\) 8 AND \(f_t\) \(\geq\) 10.

The first clause is the \emph{signal} the benchmark intends to measure.
The second is an \emph{eligibility} clause: it restricts labeling to
units that still have enough frontier left to fail against. But \(f_t\)
is monotonically non-increasing --- the frontier can only shrink as
positives are recovered or the horizon advances --- so once \(f_t\)
falls below the floor of 10, it never returns. Call a unit
\textbf{starved} at step t if \(f_t\) \textless{} 10. Starvation is
therefore \emph{absorbing}: a starved unit can never again satisfy the
eligibility clause, hence can never be labeled degraded, no matter how
badly it fails thereafter. It falls out of the ``degraded'' class and,
because the benchmark defines the control class \emph{residually}
(everything not degraded), it lands among the \textbf{controls}.

Figure 1 shows the resulting state structure. This yields a
deterministic statement that is the crux of the pathology.

\includegraphics[width=1\linewidth,height=\textheight,keepaspectratio]{fig1_absorbing_shadow_state.pdf}

\textbf{Figure 1.} \emph{The eligibility conjunct creates an absorbing
shadow state.} \textbf{(A)} The rule sends a unit to \emph{Degraded}
only on eight consecutive misses while the frontier is still eligible
(\(f_t\) \(\geq\) 10); once the frontier falls below the floor the unit
is \emph{Starved} --- an absorbing state with no exit --- and is filed,
at the horizon, as a residual control. An outcome no longer
ascertainable is not thereby a negative outcome. \textbf{(B)} Read at
the landmark (t = 20), 69.2\% of the nominal control class --- 3,292 of
4,755 --- is starved-by-the-landmark shadow class, against 1,460 clean
controls and 3 post-landmark starvers.

\textbf{Observation 1 (starvation implies control).} If a unit becomes
starved at any time before the horizon without first being degraded, it
can never subsequently become an event and therefore enters the residual
control class. This includes units already starved at the landmark and
units that starve subsequently.

An event must satisfy f \(\geq\) 10 at its degradation step. Since f is
non-increasing, once a unit is starved (f \textless{} 10) it can never
again meet the eligibility clause and can never be degraded; not being
degraded by the horizon, it is a control by the residual definition. For
a unit already starved at the landmark the argument is immediate --- it
is unlabelable at the landmark; for a unit that starves later it applies
from the starvation step onward.

The consequence is that the control class does not mean ``a healthy user
who did not degrade.'' It substantially means ``a user whose frontier
ran out --- who is failing continuously --- and whom the rule is
structurally unable to label.'' The event/control contrast the benchmark
presents is not \emph{will-degrade} versus \emph{will-not-degrade}; it
is \emph{will-degrade} versus \emph{(will-not-degrade OR
cannot-be-labeled)}.

Two terms recur below and are worth fixing precisely. The phrase
\emph{shadow class} has unrelated uses in other fields; we use it here
in a specific measurement-theoretic sense.

\textbf{Definition 1 (shadow class).} Given a population, a target
outcome, and a labeling rule, a \emph{shadow class} is a subset of units
that (i) remains relevant to the inferential contrast the labels are
used to support; (ii) is structurally prevented by the labeling
architecture from having the target outcome ascertained --- not merely
unlabeled empirically, but unlabelable by construction; and (iii) is
consequentially treated: omitted, censored, or absorbed into another
nominal class in a way that changes inference. All three conditions
matter. Eligibility rules legitimately place units out of scope; a
shadow class exists only where units relevant to the contrast become
structurally unascertainable and their residual treatment changes what
the instrument reports. The starved controls of §2.2 satisfy all three:
after starvation, degradation status is unascertainable under the rule,
and the residual definition converts that unascertainability into a
negative label. An outcome that is no longer ascertainable is not
thereby a negative outcome.

\textbf{Definition 2 (shadow-class census).} The report, prior to any
use of the benchmark for inference, of the prevalence and behaviour of
each shadow class the label rule induces.

A shadow class in this sense is not a latent class: its membership is
exactly identifiable from observables (here, \(f_t\) \textless{} 10).
Nor is it missing data. Three situations should be kept distinct:
\emph{data invisibility} --- the unit is not observed; \emph{label
uncertainty} --- the unit is observed but its correct label is unknown;
and \emph{shadow-class unascertainability} --- the unit is fully
observed, every step of its trajectory recorded, while the labeling
architecture has made one relevant outcome structurally unreachable. The
pathology reported here is the third.

\subsubsection{2.2 The empirical
signature}\label{the-empirical-signature}

The contamination is not marginal. Of the 4,755 units the rule files as
controls (every unit not labelled degraded through the horizon), read at
the landmark step (step 20):

\begin{itemize}
\tightlist
\item
  \textbf{3,292 of 4,755 controls (69.2\%) are frontier-starved by step
  20.} By Observation 1, \textbf{0 of 20,265 events are} --- verified by
  direct assertion.
\item
  The starved controls fail almost continuously: control units miss a
  \textbf{median of 15 of the 20} pre-landmark steps (up to 19). Event
  units, which by construction have not yet degraded, miss \textbf{at
  most 7 of 20} (median 1).
\end{itemize}

The resulting bias runs in one direction. The most natural early-warning
statistic --- the early miss rate --- is therefore
\textbf{anti-informative} on this benchmark: the highest-missing units
are the controls, not the events. A detector that simply reads ``who is
missing a lot early'' ranks the controls above the events. And because
any indicator computed from the recommendation stream is downstream of
the same frontier mechanics, an indicator can partially recover the
label's mechanical structure --- \emph{who is starved} --- with no
early-warning content whatsoever. The benchmark measures the wrong
quantity, and therefore presents a corrupted target to any detector
evaluated against these labels.

A separate quantity shows the importance of warning runway: of 35,584
whole-horizon events, only 20,265 remain event-free and eligible at step
20 --- \textbf{43\% had already degraded}, with no runway over which a
step-20 warning could have been assessed. These early events are
labeled, not shadow-class members; they are excluded by standard
landmark logic.

\subsubsection{2.3 A reductio via a random
recommender}\label{a-reductio-via-a-random-recommender}

The pathology is a property of the \emph{label rule's interaction with
the engine}, not of any particular recommender. Replace the
collaborative-filtering engine with a uniform-random recommender and
drive the \emph{same} frozen replay rules over the \emph{same} users:
the population and their rating histories are unchanged, and only the
engine differs. A random recommender draws a ten-item slate from the
full unseen catalogue, so frontier recovery is rare by construction:
essentially every eligible user accrues eight consecutive misses within
the first several steps and is labelled degraded immediately. Because
degradation before the landmark removes a unit from the analysis
population, the landmark population under this engine is empty ---
literally so on the pre-registered degeneracy probe, which admitted
\textbf{0 of 3,000 users}, with 63.5\% excluded for degrading before the
minimum warning runway.

That the users are held fixed is what makes this an invalidity rather
than a sensitivity. A metric may legitimately vary with the quality of
the system it measures. But the construct here claims to describe
\emph{users}, and the same users move from a fifth of the horizon
population being labelled to essentially all of them, on an engine swap
alone.

The label nominally measuring \emph{user degradation} therefore fires on
essentially everyone the moment the engine stops being a competent
recommender. That is difficult to reconcile with a construct supposed to
measure a property of users, and it is near-dispositive evidence that
the construct is substantially engine-recovery behaviour rather than a
user property.

\subsubsection{2.4 Detection: a learnability gate, evaluated before any
decision
quantity}\label{detection-a-learnability-gate-evaluated-before-any-decision-quantity}

The reason this pathology is reported here rather than discovered by a
reviewer is that the evaluation was pre-registered (OSF) with a
\textbf{learnability gate} placed upstream of the detector comparison.
The gate asks a prerequisite question: at the primary false-positive
budget, does the simplest label-derived signal --- the early miss rate
--- separate events from controls above chance? Functionally this is a
prerequisite-direction test: the outcome is defined through runs of
misses, so the simplest outcome-adjacent observable must not have its
direction structurally reversed by the negative class. On the
contaminated population it does not; its true-positive rate is zero at
the operating point, its confidence interval excluding zero on the wrong
side. Under the frozen protocol, a learnability-gate failure routes the
study to \emph{uninterpretable-by-construction}: the detector results
are reported but carry no evidential weight, because the population
cannot distinguish ``no early-warning signal'' from ``a label whose
control class is contaminated.'' The gate is evaluated \emph{before} any
detector-versus-baseline decision quantity is computed, so the finding
cannot be a post-hoc rationalization of an unwelcome result. The gate
blocked interpretation of detector results until the directional
contradiction in the benchmark had been resolved.

\subsection{Part 3 --- The contamination is structural, and how to build
a benchmark that avoids
it}\label{part-3-the-contamination-is-structural-and-how-to-build-a-benchmark-that-avoids-it}

A natural hope is that the contamination can be tuned away --- for
instance by moving the landmark earlier, so that fewer units have
starved by the time labels are read. It cannot. The following makes that
precise.

\textbf{Proposition 2 (landmark invariance).} Define a \emph{clean
control} as a unit that is neither degraded nor starved through the
horizon T. The set of clean controls is invariant to the choice of
landmark step. Consequently, no landmark choice can increase the number
of usable controls: their scarcity is a structural property of the
(floor, horizon, engine) triple, not of the evaluation design.

\emph{Proof.} Both predicates defining a clean control --- ``never
degraded by T'' and ``never starved by T'' --- are evaluated over the
full horizon and are functions of the unit's trajectory alone. The
landmark step determines only the window over which indicators are
computed and the point after which labels are read; it does not alter
any trajectory or its terminal classification. Hence the clean-control
set is identical for every landmark. Because starvation is absorbing, an
earlier landmark cannot rescue a unit that will starve later: changing
the landmark while holding T fixed cannot alter whether a trajectory
subsequently starves by T. \(\blacksquare\)

The corollary is the design lesson: \textbf{an eligibility conjunct with
a monotone, absorbing bound cannot yield a valid residual control class,
and no downstream analysis choice repairs it.} The fix must be made at
construction time. We give a specification for a benchmark of this
family that can support early-warning evaluation:

\begin{enumerate}
\def\labelenumi{\arabic{enumi}.}
\tightlist
\item
  \textbf{Decouple eligibility from labelability.} Remove the absorbing
  shadow state --- for example, a \emph{replenishing} frontier (a
  rolling window of upcoming positives) so that no unit is permanently
  barred from being labeled, and failure is always measurable against a
  live frontier.
\item
  \textbf{Define the control class positively, never residually.} A
  control should be a unit that demonstrably \emph{sustains recovery
  capability} through the horizon, not merely a unit the rule failed to
  label.
\item
  \textbf{Require a shadow-class census for any conjunctive label rule.}
  Every eligibility condition partitions the population into labelable
  and unlabelable units. Not every unlabelable stratum is pathological
  --- rules may legitimately place units out of scope --- but wherever
  the three conditions of Definition 1 co-occur, the benchmark must
  report the size and behaviour of the shadow class, because that is
  where residual contamination hides.
\item
  \textbf{Engineer prevalence at construction.} Fix a minority-class
  floor (e.g.~\(\geq\) 15\%) so that the interesting class is neither
  degenerate nor overwhelming, rather than accepting whatever prevalence
  the rule happens to produce.
\item
  \textbf{Ship a learnability gate as a required validation artifact.}
  The prerequisite check --- can the simplest label-derived signal beat
  chance? --- should be part of releasing the benchmark, not an optional
  downstream analysis.
\end{enumerate}

\textbf{Relation to prior work.} The statistical phenomenon reported
here is not new. Observation 1 describes an absorbing competing-risk
pathology, a hazard with a five-decade literature in survival analysis,
and we claim no priority over it. Starvation and degradation are
competing terminal events. The benchmark's error is to collapse the
starvation event into the residual negative class rather than represent
it as a competing outcome. When starvation is instead treated as if it
conveyed ordinary non-event status, the resulting analysis has the same
inferential pathology that motivates competing-risk methods, for which
Putter, Fiocco and Geskus (2007) is a standard tutorial and Fine and
Gray (1999) one canonical treatment. Landmarking supplies one component
of the repair by defining the at-risk population (Anderson, Cain and
Gelber, 1983; van Houwelingen, 2007). The binary benchmark additionally
excludes units that subsequently experience the competing starvation
event from the residual control class, as specified in §4.1.

Three further literatures stand in weaker relations. Reject inference is
the closest structural analogue: there, as here, labels exist only
conditional on a gate operated by the system under evaluation (Hand and
Henley, 1997; Crook and Banasik, 2004). Positive-unlabeled learning is
an inverted analogue --- its unlabeled class hides positives, whereas
ours absorbs into the negative class units for which the target outcome
has become structurally unascertainable (Elkan and Noto, 2008; Bekker
and Davis, 2020). And the eligibility-bias cluster in
pharmacoepidemiology --- depletion of susceptibles (Moride and Abenhaim,
1994), immortal-time bias, collider stratification (Suissa, 2008; Hernán
et al., 2004) --- is the direct epidemiological cousin, differing in
substrate rather than in mechanism.

One demarcation matters before it is blurred. The offline-evaluation
biases familiar in recommender research --- exposure bias and
missing-not-at-random feedback (Marlin and Zemel, 2009; Steck, 2013),
and closed-loop contamination from prior recommendation policies
(Chaney, Stewart and Engelhardt, 2018) --- are properties of how the
data were collected. The pathology here lives in the definition of the
label itself, and would persist under perfectly collected data.

Our contribution is the location of this hazard and a means of detecting
it. What appears to be unrecognized is that it is latent in
recursive-recommender degradation benchmarks whose labels are
algorithmically defined rather than observed; that it inverts the sign
of the most natural early-warning statistic, so that the contamination
is invisible in every summary the benchmark reports; and that a
pre-registered learnability gate together with a shadow-class census
detects it cheaply, before any result exists. That places the work in
the benchmark-construct-validity literature (Raji et al., 2021; Jacobs
and Wallach, 2021) rather than in survival analysis: the statistics are
borrowed, the measurement problem is the finding. The gate and the
census are what we offer for reuse.

\subsection{Part 4 --- The corrected study, as
executed}\label{part-4-the-corrected-study-as-executed}

\textbf{Statistical procedures.} Effect sizes are Cohen's \emph{d}; the
contaminated and purified columns of §4.4 report each indicator's
event-versus-control separation before and after competing-risk
purification of the control class, and the mechanical Δd for each
indicator is their difference --- the registered contrast between its
contaminated- and purified-population effect sizes. The separate
bounding diagnostic of §4.4 contrasts controls starved by the landmark
against clean controls (\textbar d\textbar{} = 5.6--6.1). Interval
estimates are bias-corrected and accelerated bootstrap intervals (BCa;
Efron, 1987), resampled over the clustering unit wherever a statistic
pools across units. Detectors are compared at matched control-alarm
counts rather than at fixed score thresholds --- k = \{15, 38, 77, 153\}
at the four registered budgets --- so every comparison sits at an
identical false-positive count, with the operating point taken as the
largest matched alarm count not exceeding the budget. The
z-standardisation parameters for all indicators are estimated once on
the 694 calibration-fold clean controls and applied unchanged out of
fold.

\subsubsection{4.1 What was corrected, and what was
audited}\label{what-was-corrected-and-what-was-audited}

Study 1's benchmark could not support the measurement it was built for.
Study 2 was pre-registered (OSF, osf.io/wka72) as its repair: starvation
is treated as a competing terminal event, the control class is
restricted to units that neither degrade nor starve through the horizon,
and every downstream quantity --- indicator normalisation, matched alarm
budgets, covariate adjustment --- is computed against that purified
population. The analysis plan, the four acceptance criteria, and the
response to each possible verdict were all fixed before the
clean-control census existed.

One methodological point should be stated exactly. Landmarking defines
the at-risk population at L; the repair here additionally treats
subsequent starvation as a competing terminal event and excludes that
stratum from the residual control class, which goes beyond landmark
analysis proper. The repaired benchmark's estimand is therefore: among
units event-free and eligible at the landmark, does a detector
distinguish subsequent degradation from sustained
non-degradation-and-non-starvation through the horizon?

Between registration and execution, an exhaustive conformance audit of
the executed configuration against the registration returned thirty of
thirty-four parameters conformant and four departures: three arithmetic
divergences in registered quantities --- the analysis population filter,
the matched-alarm budget grid, and a pre-census table inheriting the
same denominator error --- and one labelling defect, since corrected.
All four share one root cause: configuration constants transcribed from
the Study 1 registration rather than the operative Study 2 one. A
parallel sweep found three artifact writers absent from the versioned
tree; each was committed, and its artifact content-verified or
canonically re-derived to byte identity. Three further findings,
detailed in the supplement, concerned input provenance and the
completeness of a process attestation.

Each finding was recorded in the deviations log before it was
investigated. The provenance chain is content-pinned end to end, from
the published source-dataset hash through the derived panel, the
clean-control census and the covariate cache, to the verdict artifact.

Two of the divergences invalidated an earlier execution, which ran on a
population appearing in no registered analysis. \textbf{That execution
had halted at its own learnability gate and produced no detector
comparison; vacating it overturned no result.} Its record is retained
and its verdict vacated. \textbf{The analysis reported below is the
first execution of the registered design.}

\subsubsection{4.2 The repaired instrument
measures}\label{the-repaired-instrument-measures}

Before any detector comparison, the pre-registered protocol evaluates a
learnability gate: on the analysis population, does the most direct
label-derived signal --- the early miss rate --- separate events from
controls above the matched-budget chance rate? On Study 1's contaminated
population it did not, which is what routed that study to
\emph{uninterpretable}.

On the purified population it does. The gate passes with a true-positive
rate of \textbf{0.561, BCa 95\% CI {[}0.549, 0.574{]}} at the primary
budget, against a chance rate of 0.05 --- an excess over chance of 0.511
{[}0.499, 0.524{]}.

We report this as a primary result. It establishes that competing-risk
purification, applied to a benchmark exhibiting the shadow-class
pathology, recovers a population that supports the registered binary
early-warning contrast --- a single-instance demonstration that the
repair works, on the benchmark where the pathology was found.

The remaining gates pass likewise: the liveness check on indicator
distributions, and the power gate at 10,153 evaluation-fold events
against a floor of 200.

Analysis population: 20,265 events and 1,460 clean controls, split 694
calibration and 766 evaluation by the pre-registered seedless hash rule.
Matched alarm counts at the registered budgets \{0.02, 0.05, 0.10,
0.20\} are k = \{15, 38, 77, 153\}.

\subsubsection{4.3 Verdict: the signal
dies}\label{verdict-the-signal-dies}

The pre-registered composite is S = z(G) + z(p) − z(\(\delta\)),
evaluated out of fold at matched control-alarm counts against six
baseline families and the early-miss-rate reference. Three of the four
acceptance criteria are evaluated and all three fail; the fourth, which
concerns mechanism specificity of a detected signal, is not reached.

\textbf{Criterion 1 --- discrimination advantage at the primary budget.}
Required: ΔTPR \(\geq\) 0.05 with a cluster-aware BCa interval excluding
zero. Observed: \textbf{ΔTPR = −0.472, CI {[}−0.496, −0.444{]}}. The
composite does not fall short of the bar; it is substantially worse than
the best baseline, with the entire interval on the opposite side of
zero.

\textbf{Criterion 2 --- sign consistency across budgets.} Required:
positive at three of four. Observed: negative at four of four ---
\textbf{−0.503, −0.472, −0.436, −0.386} at b = \{0.02, 0.05, 0.10,
0.20\}. The disadvantage is monotone in budget and never approaches
parity.

\textbf{Criterion 3 --- covariate adjustment.} Required: coefficient
signs in the hypothesised directions (bG \textgreater{} 0, bp
\textgreater{} 0, b\(\delta\) \textless{} 0) and incremental AUC
\(\geq\) 0.02 over the covariates alone. Observed:

{\def\LTcaptype{none} % do not increment counter
\begin{longtable}[]{@{}
  >{\centering\arraybackslash}p{(\linewidth - 8\tabcolsep) * \real{0.2000}}
  >{\centering\arraybackslash}p{(\linewidth - 8\tabcolsep) * \real{0.2000}}
  >{\centering\arraybackslash}p{(\linewidth - 8\tabcolsep) * \real{0.2000}}
  >{\centering\arraybackslash}p{(\linewidth - 8\tabcolsep) * \real{0.2000}}
  >{\centering\arraybackslash}p{(\linewidth - 8\tabcolsep) * \real{0.2000}}@{}}
\toprule\noalign{}
\endhead
\bottomrule\noalign{}
\endlastfoot
\textbf{coefficient} & \textbf{estimate} & \textbf{BCa 95\% CI} &
\textbf{required} & \textbf{met} \\
bG (churn) & −0.475 & {[}−0.635, −0.313{]} & \textgreater{} 0 & no \\
bp (occupancy) & −0.419 & {[}−0.512, −0.328{]} & \textgreater{} 0 &
no \\
b\(\delta\) (coverage) & −0.538 & {[}−0.691, −0.385{]} & \textless{} 0 &
\textbf{yes} \\
\end{longtable}
}

Incremental AUC = \textbf{0.162, CI {[}0.143, 0.183{]}} --- eight times
the required floor. The criterion fails on the two inverted
coefficients.

Two things follow. The indicators carry considerable discriminative
information beyond the covariates, and for two of the three the
relationship runs opposite to the hypothesis. But the coverage
coefficient satisfies its registered direction with an interval clear of
zero. The criterion's failure is not uniform, and reporting it as
uniform would misdescribe the result.

Two frames must be kept distinct when counting. In the \textbf{joint
model} two coefficients are inverted and one is consistent. On
\textbf{marginal effects} (§4.4) one indicator is inverted, one is null
and one is consistent. Occupancy illustrates why the frames differ: its
marginal effect is +0.05, effectively null, while its partial
coefficient is −0.419 with an interval clear of zero. Partial and
marginal effects can legitimately diverge under shared variance, and
that divergence is expected here rather than anomalous.

The predictors are correlated but not severely so: on the
standardisation population the variance inflation factors are
approximately 3.5 for churn, 3.3 for coverage and 1.3 for occupancy ---
moderate by conventional standards, and not high enough to render the
estimates unstable. We therefore let the coefficients stand as reported.
What this implies for the theory: bG at −0.475 with a clean interval and
a moderate VIF is a robust inverted partial effect, which is more
damaging to the hypothesis than an unstable one would be.

\textbf{We do not fit a re-signed composite.} Inverted coefficients
invite the obvious follow-up --- a composite with the signs corrected to
match the data. Under the registration that is a post-hoc hypothesis
constructed after seeing the result, and it is foreclosed. We publish
the incremental AUC of 0.162 as the constraint any future, separately
registered hypothesis would have to respect, and pursue nothing further
here.

In absolute terms the composite reaches a matched true-positive rate of
\textbf{0.089} at the primary operating point --- against a chance rate
of 0.05, and against \textbf{0.561} for the best baseline. Figure 3
shows the discrimination gap at all four budgets.

\pandocbounded{\includegraphics[keepaspectratio]{fig3_delta_tpr_by_budget.pdf}}

\textbf{Figure 3.} \emph{Discrimination advantage at matched alarm
budgets.} Composite minus best baseline (ΔTPR) at the four registered
budgets, with cluster-aware BCa 95\% intervals; the parity (ΔTPR = 0)
and acceptance (ΔTPR \(\geq\) 0.05) lines are marked. Every point is
negative with its interval clear of zero --- the composite is strictly
worse than the best miss-based baseline at every budget.

The baseline set is narrower than it looks. Five of the seven
comparators --- the early-miss reference, variance EWS, CUSUM,
Page--Hinkley and matrix profile --- achieve \emph{identical} matched
true-positive rates of 0.561. All are computed from the same derived
miss-run series, and at matched alarm counts they select the same units.
They are therefore not seven independent baselines but essentially two:
a miss-based family and permutation entropy, the latter below chance at
0.041 --- under the 0.05 matched-budget chance rate here, and under the
matched chance rate at all four budgets (0.016, 0.041, 0.093, 0.183
against 0.02, 0.05, 0.10, 0.20), so its sub-chance behaviour is
consistent across budgets rather than a single-budget artifact. This
weakens the ``best of six families'' framing while leaving the
comparison itself intact, since the composite loses to the miss-based
family by a wide margin however that family is counted.

\subsubsection{4.4 Characterisation}\label{characterisation}

Six pre-registered analyses and one post-verdict descriptive comparison
characterise the result. The scope of the six was fixed and anchored
before the sealed decision quantities were opened. The seventh was
granted afterwards under the same default-denied mechanism, with its
interpretive rule and reporting threshold fixed and anchored before its
values existed; it is identified as such below.

\textbf{The apparent Study 1 signal was contamination.} All four
indicators show large effects on the contaminated population.
Purification removes or reverses every one of them (Figure 2):

{\def\LTcaptype{none} % do not increment counter
\begin{longtable}[]{@{}
  >{\centering\arraybackslash}p{(\linewidth - 6\tabcolsep) * \real{0.2500}}
  >{\centering\arraybackslash}p{(\linewidth - 6\tabcolsep) * \real{0.2500}}
  >{\centering\arraybackslash}p{(\linewidth - 6\tabcolsep) * \real{0.2500}}
  >{\centering\arraybackslash}p{(\linewidth - 6\tabcolsep) * \real{0.2500}}@{}}
\toprule\noalign{}
\endhead
\bottomrule\noalign{}
\endlastfoot
\textbf{indicator} & \textbf{contaminated} & \textbf{purified} &
\textbf{mechanical Δd {[}95\% CI{]}} \\
G (churn) & +2.01 & \textbf{−0.64} & 2.65 {[}2.55, 2.74{]} \\
p (occupancy) & −2.25 & \textbf{+0.05} & −2.30 {[}−2.39, −2.20{]} \\
\(\delta\) (coverage) & +2.00 & \textbf{−0.68} & 2.68 {[}2.58,
2.78{]} \\
miss-run fraction & −2.62 & \textbf{+0.81} & −3.43 {[}−3.55, −3.32{]} \\
\end{longtable}
}

\pandocbounded{\includegraphics[keepaspectratio]{fig2_bounding_diagnostic.pdf}}

\textbf{Figure 2.} \emph{Every large contaminated-population effect is
dominated by the mechanical contamination.} Paired Cohen's d for all
four indicators on the contaminated population (vermillion circles, vs
all controls) and the purified population (ink squares, vs clean
controls); each large contaminated effect collapses toward zero or
reverses under purification.

Two distinct frames must be kept apart here, and conflating them would
misstate the result. Relative to the \emph{contaminated} values, all
four quantities change substantially: G and \(\delta\) invert, p
collapses, and the miss-run fraction reverses. Relative to the
\emph{hypothesis}, the picture is not uniform. The theory predicted G
high, p high and \(\delta\) low in events. On the purified population G
is inverted (−0.64), p is null (+0.05), and \textbf{\(\delta\) is in the
predicted direction (−0.68)}. The miss-run fraction's reversal runs the
other way again --- from anti-informative on the contaminated population
to informative on the purified one, the same flip the learnability gate
shows.

So one of the three registered slate indicators retains a moderate
effect consistent with the hypothesis.

The mechanical component of the shift has a cluster-bootstrap interval
excluding zero for all four quantities. The contrast against the starved
stratum identifies what the contaminated effects were tracking. In this
bounding diagnostic --- evaluation-fold starved controls against clean
controls --- the slate indicators reach \textbar d\textbar{} of 5.6 to
5.7 and the miss-run fraction 6.1. (This is a distinct comparison from
the pre-landmark probe reported next, and its d range should not be
merged with that one.)

\textbf{And the mechanism is visible directly.} In the pre-landmark
starvation-concurrence probe --- controls starved by the landmark (n =
3,292) against clean controls (n = 1,460), indicators taken over steps
1--20 --- the composite separates them at \textbf{d = +4.20}, with
individual indicators from −5.77 to +6.29. The composite is, in
substantial part, a frontier-depletion detector. In the contaminated
population that is exactly what it was rewarded for detecting;
purification removes the thing it was reading.

\textbf{The null is robust to landmark choice.} Repeating the full
primary analysis at landmarks 15 and 25: criteria 1 and 2 fail at every
landmark; criterion 3 fails at 20 and 25 and is coverage-limited at 15
(58.65\% covariate coverage, declared before computation).

\textbf{The indicators are substantially correlated.} On the
standardisation population --- the 694 calibration-fold clean controls
from which the z-parameters derive --- the pairwise correlations are
r(G, \(\delta\)) = \textbf{0.833}, r(G, p) = \textbf{−0.478}, r(p,
\(\delta\)) = \textbf{−0.415}. The composite's three terms are therefore
not three independent axes, though neither are they a single quantity
measured three times.

The population matters here. Computed instead on the full
evaluation-fold arm, which includes events and starved controls, the
same correlations are far higher (0.975, −0.869, −0.869). That inflation
is itself attributable to the contamination: the starved stratum carries
extreme indicator values, and its presence drives the indicators into
apparent lockstep. Label contamination changed not only estimated
discrimination but the apparent covariance structure of the proposed
measurement instrument. The values reported above are those on the
population where standardisation is defined.

The generalisable point follows. \textbf{Feature geometry estimated on a
population whose labels have not been validated can inherit the
contamination.} A composite, index, factor structure or weighting scheme
designed from such correlations can therefore be designed partly around
an artifact --- here, a difference between 0.975 and 0.833 that would
materially change how one reasons about the feature space. Taken with
the label itself and the effect sizes, the same contamination has now
been shown to corrupt three distinct layers of the analysis.

\textbf{The composite carries nothing on severed arms.} Absolute
discrimination against chance, reported unconditionally under the
pre-registered protocol, exceeds chance on \textbf{zero of six} control
arms in which the recursive personalisation mechanism is disabled. Four
arms are significantly below chance; on two of the shuffled-timestamp
arms the interval includes zero, so those are indistinguishable from
chance rather than below it.

\textbf{Frontier headroom does not localise the null.} Events with
remaining frontier above the eligibility threshold reproduce the primary
result (\textbf{ΔTPR = −0.470}); the narrow-headroom stratum contains 38
events and is reported as underpowered rather than estimated.

\textbf{Who the label discards} \emph{(the post-verdict comparison)}.
Starvation is absorbing and never co-occurs with an event: \textbf{zero
of 20,265 events} cross the eligibility floor before collapse,
confirming Observation 1 empirically. The discarded stratum is not a
random sample of users --- those excluded have thinner early positive
histories than clean controls (\textbf{d = −0.441, CI {[}−0.500,
−0.381{]}}), a moderate association. This comparison was granted after
the verdict, with its direction, effect-size threshold and reporting
rule fixed and anchored before any value was computed. The measured
variable is warm-start positive count, not total activity; no claim
about usage frequency is made or licensed, and no other covariate was
compared.

Three registered analyses were not run and are reported with their
reasons: the matrix-factorisation and sequential control engines existed
to establish mechanism specificity of a detected signal, and none was
detected; the random-engine re-run is moot, its degeneracy being the
result itself; and the starvation-discrimination diagnostic as
originally framed is vacuous by population fact, with three
within-window starvers against 1,460 clean controls.

\subsection{Part 5 --- What this shows}\label{part-5-what-this-shows}

\subsubsection{5.1 Conflation, separation, and what remains
unexplained}\label{conflation-separation-and-what-remains-unexplained}

\textbf{The label conflated two signals.} One was mechanical --- the
interaction between the recommender and a depleting frontier, which the
composite detects at d = +4.20 and which the eligibility conjunct
silently promoted into the control class. The other was a genuine
early-warning signal in the outcome stream, which becomes informative
once purification removes the contamination and which drives the
learnability gate. In the contaminated benchmark these were welded
together, and the apparent success of the slate composite was almost
entirely the first.

\textbf{A first-order consistency check explains the near-cancellation:}
Combining the purified marginals as the registered composite combines
them --- z(G) + z(p) − z(\(\delta\)) --- gives (−0.64) + (+0.05) −
(−0.68) \(\approx\) +0.09, against a composite standard deviation of
\(\sqrt{3 + 2[r(G,p) - r(G,\delta) - r(p,\delta)]} = \sqrt{1.208} \approx 1.10\).
The inverted churn term approximately cancels the coverage term it is
subtracted from, the null occupancy term carries the remainder, and the
net standardised separation is roughly 0.08. That is near chance, and it
is consistent with the observed matched true-positive rate of 0.089
against a chance rate of 0.05. The figure is approximate: it treats the
marginal effect sizes as though pooled-SD standardised.

So the composite's failure is arithmetically consistent with its own
components. \textbf{What we do not explain is why churn inverts on the
purified population at all.} That is the deeper question, and this
design cannot answer it; we decline to speculate.

What the result establishes is narrower. On this benchmark, at this
landmark, the registered composite carries no usable early-warning
signal; on marginal effects one component points in the predicted
direction, one is null and one is inverted; and the outcome-stream
quantity the design treated as a baseline outperforms all of them.

The single sentence: \textbf{On this benchmark, the outcome-stream
baseline carries substantially more early-warning information than the
registered slate-structure composite --- and the theory placed its
signal on the wrong side of the instrument.}

\subsubsection{5.2 What the repair validated, and what remains
proposed}\label{what-the-repair-validated-and-what-remains-proposed}

Three elements of the §3 specification were exercised in this case:
competing-event purification restored the intended binary contrast; the
shadow-class census exposed the magnitude and behaviour of the excluded
stratum; and the pre-registered gate confirmed that the repaired
population recovers the required directional prerequisite. Together,
these three operations exposed the contamination, defined the repaired
contrast, and confirmed that the registered directional prerequisite was
restored.

Two elements remain proposed rather than validated: the replenishing
frontier and engineered minority-class prevalence. Their necessity is
partially evidenced --- the invariance result shows that no landmark
choice rescues a construction with an absorbing eligibility bound, which
is the argument for decoupling eligibility from labelability at
construction time. Their sufficiency is untested here, because the
repair adopted was purification of an existing benchmark rather than
construction of a new one.

The distinction matters for anyone adopting the specification: the
diagnostic half was exercised and demonstrated in this instance, the
constructive half is proposed.

\subsubsection{5.3 What the negative does and does not
establish}\label{what-the-negative-does-and-does-not-establish}

The pre-registered composite showed no advantage over miss-based
baselines at matched alarm budgets, no consistent sign across budgets,
and no incremental discrimination in the hypothesised direction --- on
this benchmark, this engine family, this horizon, and this landmark,
robust across two alternative landmarks.

That is decisive against the registered hypothesis. It is not a
demonstration that no slate-derived statistic could carry early-warning
information in any recursive recommender. Three limits on the claim:

The composite was pre-specified, not searched. Its three terms proved to
be substantially correlated rather than independent; a different
construction is not tested here and is not excluded by this result.

The population is definitionally restricted. Units that starve before
the landmark are unlabelable by construction in this benchmark family
--- nothing about the composite was measured on them and nothing could
be, under any design in this family. This is a scope statement, not a
qualification of the verdict, and it is one the specification already
makes.

The setting is one offline replay, one item-item collaborative-filtering
engine family, one dataset. Cross-domain generalisation was never tested
and is not claimed.

\subsubsection{5.4 Implications for benchmark
construction}\label{implications-for-benchmark-construction}

\textbf{The prescriptions below apply to algorithmically labeled
benchmarks in which an eligibility condition can render the target
outcome structurally unascertainable while the residual class remains
inferentially active.} Any conjunctive label rule containing an
eligibility condition should be audited for that possibility: audit
broadly, claim pathology narrowly.

First, any such evaluation should publish a shadow-class census: the
size and behaviour of the population the rule renders unlabelable. In
the case reported here that population comprised 69\% of the nominal
control class and missed a median 15 of 20 pre-landmark steps, while
events missed at most 7. No summary statistic the benchmark reported
revealed it.

Second, label validity should be established \emph{before} the feature
geometry used to design an instrument is estimated. Correlations, factor
structures and composite weightings computed on an unvalidated-label
population inherit whatever contamination that population carries; the
difference reported in §4.4 between 0.975 and 0.833 is a demonstration
of the size such an artifact can reach. Designing a detector on the
first figure would have produced a different and wrong understanding of
the feature space: once the class ontology is wrong, feature selection
and composite construction performed downstream inherit that error
before any detector is evaluated.

Third, a learnability gate should be a release artifact rather than an
optional downstream analysis. The check is cheap --- can the most direct
label-derived signal beat chance on this population? --- and its failure
flags that a null detector result cannot safely be interpreted as
absence of early-warning information until label validity is
investigated. In this project the same gate produced both answers on the
same benchmark, before and after repair.

A more general observation follows from the audit rather than the
experiment. Conformance between a registration and its executed
configuration is not established by the existence of a registration.
Three of the four departures found here were arithmetic divergences in
registered quantities, all traceable to constants transcribed from the
wrong document; none would have been visible in any reported result.
Machine-checkable registration parameters, asserted against
configuration at run time, would have caught all three.

\subsubsection{5.5 The status of the
question}\label{the-status-of-the-question}

The hypothesis that slate-structure statistics carry early warning of
recursive degradation has been tested once, on a valid instrument, and
failed on its own pre-registered terms.

The broader question --- whether early warning is achievable in
recursive recommender systems --- is not answered by this work.
Answering it requires a benchmark built to the constructive half of the
specification, or a live system, and the analysis reported here places
the constraint precisely: the present results point future work toward
the outcome stream and toward candidate quantities whose independence
from frontier depletion can be demonstrated.

A further extension, not attempted here, is a taxonomy of shadow classes
across label architectures: the absorbing variety studied here is one
type; the gated labels of reject inference suggest another. The census
of Definition 2 generalises accordingly.

\subsubsection{5.6 Limitations}\label{limitations}

The benchmark is a deterministic offline replay, not a deployed system
with feedback. The engine family is item-item collaborative filtering;
the two mechanism-control engines that would have characterised
alternative families were not run, because their registered purpose was
to establish specificity of a detected signal and none was detected. The
covariate set is fixed at two pre-registered variables.

The baseline set is less independent than its nomenclature suggests.
Five of the seven comparators are computed from the same derived
miss-run series and return identical matched true-positive rates; only
one comparator is genuinely distinct, and it performs below chance. A
stronger evaluation would draw baselines from structurally different
signal families rather than from transformations of a single series. The
verdict does not turn on this: the composite's \emph{absolute} matched
true-positive rate is 0.089 against a chance rate of 0.05, so no
baseline set of any diversity would have rescued it. Effect sizes on the
discarded stratum are descriptive and were computed under a scope
granted after the verdict, with the interpretive rule fixed before the
values existed.

The instrument now measures. What it measures is one narrow question,
and the answer is no.

\subsection{References}\label{references}
\sloppy

\emph{All entries verified against source (2026-07-22; additions
verified 2026-08-14).}

Anderson, J. R., Cain, K. C., and Gelber, R. D. (1983). Analysis of
survival by tumor response. \emph{Journal of Clinical Oncology} 1(11),
710--719.

Bekker, J. and Davis, J. (2020). Learning from positive and unlabeled
data: a survey. \emph{Machine Learning} 109(4), 719--760.
doi:10.1007/s10994-020-05877-5

Boettiger, C. and Hastings, A. (2012). Quantifying limits to detection
of early warning for critical transitions. \emph{Journal of the Royal
Society Interface} 9(75), 2527--2539. doi:10.1098/rsif.2012.0125

Chaney, A. J. B., Stewart, B. M., and Engelhardt, B. E. (2018). How
algorithmic confounding in recommendation systems increases homogeneity
and decreases utility. \emph{RecSys '18}.

Crook, J. and Banasik, J. (2004). Does reject inference really improve
the performance of application scoring models? \emph{Journal of Banking
\& Finance} 28(4), 857--874. doi:10.1016/j.jbankfin.2003.10.010

Efron, B. (1987). Better bootstrap confidence intervals. \emph{Journal
of the American Statistical Association} 82(397), 171--185.
doi:10.1080/01621459.1987.10478410

Elkan, C. and Noto, K. (2008). Learning classifiers from only positive
and unlabeled data. In \emph{Proceedings of the 14th ACM SIGKDD
International Conference on Knowledge Discovery and Data Mining (KDD
'08)}, 213--220.

Fine, J. P. and Gray, R. J. (1999). A proportional hazards model for the
subdistribution of a competing risk. \emph{Journal of the American
Statistical Association} 94(446), 496--509.

Hand, D. J. and Henley, W. E. (1997). Statistical classification methods
in consumer credit scoring: a review. \emph{Journal of the Royal
Statistical Society: Series A} 160(3), 523--541.

Harper, F. M. and Konstan, J. A. (2015). The MovieLens datasets: history
and context. \emph{ACM Transactions on Interactive Intelligent Systems}
5(4), Article 19. doi:10.1145/2827872

Hernán, M. A., Hernández-Díaz, S., and Robins, J. M. (2004). A
structural approach to selection bias. \emph{Epidemiology} 15(5),
615--625.

Jacobs, A. Z. and Wallach, H. (2021). Measurement and fairness. In
\emph{Proceedings of the 2021 ACM Conference on Fairness,
Accountability, and Transparency (FAccT '21)}, 375--385.
doi:10.1145/3442188.3445901

Marlin, B. M. and Zemel, R. S. (2009). Collaborative prediction and
ranking with non-random missing data. \emph{RecSys '09}.

Moride, Y. and Abenhaim, L. (1994). Evidence of the depletion of
susceptibles effect in non-experimental pharmacoepidemiologic research.
\emph{Journal of Clinical Epidemiology} 47(7), 731--737.
doi:10.1016/0895-4356(94)90170-8

Putter, H., Fiocco, M., and Geskus, R. B. (2007). Tutorial in
biostatistics: competing risks and multi-state models. \emph{Statistics
in Medicine} 26(11), 2389--2430. doi:10.1002/sim.2712

Raji, I. D., Bender, E. M., Paullada, A., Denton, E., and Hanna, A.
(2021). AI and the everything in the whole wide world benchmark. In
\emph{NeurIPS Track on Datasets and Benchmarks}. arXiv:2111.15366

Steck, H. (2013). Evaluation of recommendations: rating-prediction and
ranking. \emph{RecSys '13}.

Suissa, S. (2008). Immortal time bias in pharmacoepidemiology.
\emph{American Journal of Epidemiology} 167(4), 492--499.
doi:10.1093/aje/kwm324

van Houwelingen, H. C. (2007). Dynamic prediction by landmarking in
event history analysis. \emph{Scandinavian Journal of Statistics} 34(1),
70--85. doi:10.1111/j.1467-9469.2006.00529.x

\textbf{Citation integrity.} Every listed reference carries a specific
in-body citation, and every author--year citation in the text resolves
to a listed entry --- including the §3 pharmacoepidemiology attributions
(Suissa, 2008; Hernán et al., 2004). No orphans in either direction.

\subsection{Data, code and pre-registration
availability}\label{data-code-and-pre-registration-availability}

Study 1 pre-registration (frozen document, OSF project osf.io/7bvgz).
Study 2 pre-registration: OSF osf.io/wka72. Source data: MovieLens-25M,
GroupLens Research, verified against the published distribution hash.
Analysis code, frozen configuration, complete deviations log, and hash
manifests: https://github.com/davidmullett/loopzero-paper-public.
Derived panels are embargoed pending release; their content hashes are
published in the input manifest and released artifacts verify against
them.

\subsection{Competing interests}\label{competing-interests}

The author is the founder of Loopzero, Inc., which has filed a U.S.
provisional patent application (No.~64/042,661, filed April 17, 2026)
covering monitoring and verification methods related to the subject
matter of this work; the author is the named inventor. The author
declares no other competing interests.

\end{document}
